\subsection{ENVELOPING BIGEBRA OF A LIE ALGEBRA}

Recall that g denotes a Lie algebra and U its enveloping algebra, with its canonical filtration \((\mathrm{U}_{n})_{n\geqslant0}\) .

PROPOSITION 7. There exists on the algebra U one and only one coproduct c which makes U into a bigebra such that the elements of \(\sigma(g)\) are primitive. The bigebra (U, c) is cocommutative; its counit is the linear form \(\varepsilon\) such that the constant term (Chapter I, §2, no. 1) of every element \(x\) of U is \(\varepsilon(x)\) .1. The canonical filtration \((\mathrm{U}_n)_{n\geqslant 0}\) of U is compatible with this bigebra structure.

\begin{enumerate}
\def\labelenumi{(\alph{enumi})}
\tightlist
\item
  Let \(x \in \mathfrak{g}\) ; we write \(c_0(x) = \sigma(x) \otimes 1 + 1 \otimes \sigma(x) \in \mathrm{U} \otimes \mathrm{U}\) . If \(x, y\) are in \(\mathfrak{g}\) , then
\end{enumerate}

\[
c _ {0} (x) c _ {0} (y) = (\sigma (x) \sigma (y)) \otimes 1 + 1 \otimes (\sigma (x) \sigma (y)) + \sigma (x) \otimes \sigma (y) + \sigma (y) \otimes \sigma (x),
\]

whence

\[
[ c _ {0} (x), c _ {0} (y) ] = c _ {0} ([ x, y ]).
\]

By the universal property of U (Chapter I, §2, no. 1, Proposition 1), there exists one and only one unital algebra homomorphism

\[
c \colon \mathrm{U} \to \mathrm{U} \otimes \mathrm{U}
\]

such that \(c(\sigma(x)) = \sigma(x) \otimes 1 + 1 \otimes \sigma(x)\) for \(x \in \mathfrak{g}\) . This proves the uniqueness assertion of Proposition 7.

\begin{enumerate}
\def\labelenumi{(\alph{enumi})}
\setcounter{enumi}{1}
\tightlist
\item
  We show that \(c\) is coassociative. The linear mappings \(c'\) and \(c''\) of \(U\) into \(U \otimes U \otimes U\) defined by
\end{enumerate}

\[
c ^ {\prime} = (c \otimes \mathrm{Id} _ {\mathrm{U}}) \circ c \quad \text { and } \quad c ^ {\prime \prime} = (\mathrm{Id} _ {\mathrm{U}} \otimes c) \circ c
\]

are unital algebra homomorphisms which coincide on \(\sigma (\mathfrak{g})\) since, for \(a\in \sigma (\mathfrak{g})\)

\[
c ^ {\prime} (a) = a \otimes 1 \otimes 1 + 1 \otimes a \otimes 1 + 1 \otimes 1 \otimes a = c ^ {\prime \prime} (a),
\]

whence the result.

\begin{enumerate}
\def\labelenumi{(\alph{enumi})}
\setcounter{enumi}{2}
\item
  We show that \(c\) is cocommutative. Let \(\tau\) be the automorphism of \(\mathbf{U} \otimes \mathbf{U}\) such that \(\tau(a \otimes b) = b \otimes a\) for \(a, b\) in \(\mathbf{U}\) . The mappings \(\tau \circ c\) and \(c\) of \(\mathbf{U}\) into \(\mathbf{U} \otimes \mathbf{U}\) are unital algebra homomorphisms which coincide on \(\sigma(\mathfrak{g})\) , whence the result.
\item
  We show that \(\varepsilon\) is a counit for \(c\) . The mappings \((\mathrm{Id}_{\mathbb{U}} \otimes \varepsilon) \circ c\) and \((\varepsilon \otimes \mathrm{Id}_{\mathbb{U}}) \circ c\) , of U into U are unital algebra homomorphisms which coincide with \(\mathrm{Id}_{\mathbb{U}}\) on \(\sigma(\mathfrak{g})\) .
\item
  We know that \(\mathrm{U}_0 = \mathrm{K}.1, \mathrm{U}_n \subset \mathrm{U}_{n+1}, \mathrm{U} = \bigcup_{n \geqslant 0} \mathrm{U}_n\) and \(\mathrm{U}_n. \mathrm{U}_m \subset \mathrm{U}_{n+m}\) (Chapter I, §2, no. 6). Let \(a_1, \ldots, a_n\) be in \(\sigma(\mathfrak{g})\) . Then
\end{enumerate}

\[
\begin{array}{l} c (a _ {1} \ldots a _ {n}) = \prod_ {i = 1} ^ {n} c (a _ {i}) = \prod_ {i = 1} ^ {n} (a _ {i} \otimes 1 + 1 \otimes a _ {i}) \\ = \sum_ {i = 0} ^ {n} \sum_ {\alpha \in I (i)} (a _ {\alpha (1)} \ldots a _ {\alpha (i)}) \otimes (a _ {\alpha (i + 1)} \ldots a _ {\alpha (n)}), \end{array}\tag{8}
\]

where \(I(i)\) denotes the set of permutations of \([1, n]\) which are increasing in each of the intervals \([1, i]\) and \([i + 1, n]\) . As \(U_{n}\) is the K-module generated by the products of at most n elements of \(\sigma(g)\) , formula (8) implies that the filtration \((\mathrm{U}_{n})\) is compatible with the bigebra structure of \((\mathrm{U}, c)\) .

DEFINITION 3. The bigebra \((\mathrm{U}, c)\) is called the enveloping bigebra of the Lie algebra \(\mathfrak{g}\) .

PROPOSITION 8. Let E be a bigebra with coproduct denoted by \(c_{\mathbb{E}}\) and let h be a Lie algebra homomorphism of g into P(E) (no. 2, Proposition 4). The unital algebra homomorphism f: U → E such that f(σ(x)) = h(x) for all x ∈ g is a bigebra morphism.

\[
(f \otimes f) \circ c = c _ {\mathrm{E}} \circ f.
\]

\[
a \in \sigma (\mathfrak {g})
\]

\[
(f \otimes f) (c (a)) = f (a) \otimes 1 + 1 \otimes f (a) = c _ {\mathbb {E}} (f (a))
\]

since \(f(a) \in \mathrm{P}(\mathrm{E})\) . Similarly if \(\varepsilon_{\mathbb{E}}\) is the counit of \(\mathrm{E}\) , \(\varepsilon_{\mathbb{E}} \circ f\) is a unital algebra homomorphism \(\mathrm{U} \to \mathrm{K}\) which is zero on \(\sigma(\mathfrak{g})\) (no. 1, Proposition 1) and therefore coincides with \(\varepsilon\) .

It follows from Propositions 5 and 8 that the mapping \(f \mapsto f \circ \sigma\) defines a one-to-one correspondence between bigebra homomorphisms \(U(\mathfrak{g}) \to E\) and Lie algebra homomorphisms \(\mathfrak{g} \to P(E)\) .

COROLLARY. Let \(\mathfrak{g}_i\) ( \(i = 1,2\) ) be a Lie algebra, \(\mathrm{U}(\mathfrak{g}_i)\) its enveloping bigebra and \(\sigma_i\colon \mathfrak{g}_i\to \mathrm{U}(\mathfrak{g}_i)\) the canonical mapping. For every Lie algebra homomorphism \(h\colon \mathfrak{g}_1\to \mathfrak{g}_2\) , the unital algebra homomorphism \(\mathrm{U}(h)\colon \mathrm{U}(\mathfrak{g}_1)\to \mathrm{U}(\mathfrak{g}_2)\) such that \(\mathrm{U}(h)\circ \sigma_1 = \sigma_2\circ h\) (Chapter I, § 2, no. 1) is a bigebra morphism.

